\section{Prominent Attorney Marries Actor}

Ethics is the science of determining which ends man should achieve and technics is science of creating the means to achieve those ends. People act to achieve some good, usually without any deep consideration about those goods. Contemporary thought, both religious and non-religious, has relegated ethics to a corner of human life.

However, if metaphysics answers the question ``who am I?'', then ethics supplies the answer to, ``who should I become?'' The general answer is that the good is to actualize one's own nature. Here ethics is understood as all-encompassing and the opposite of mere moralizing. In other words, man is a compound of actuality and potential, and the task is to actualize all his potentials. Most people have a basic idea of this, for example, as revealed in the phrase ``be all you can be.'' People will say of someone that he is not living up to his potential, and so on.

A fundamental metaphysical principle is that only something already real can actualize a potential. In man's case, the Real I, as the intelligence of the being, can bring about that process via his true will. The goodness, then, of a man or a woman consists in how close he or she actualizes the Absolute Man or Absolute Woman.

Unfortunately, the mass of men have no means to judge that. That is why sports are so popular. There is no debate about who the best golfers or footballers are. Similarly, in entertainment, we know who the most attractive personalities are. Other easily judged categories are wealth, power, and finally popularity, even if one is popular just for being popular. So, it is easy to conclude that one should actualize athletic prowess, beauty, power, or riches.

Clearly, these cannot represent the meaning of life, since they depend on external factors, not from one's own specific being. But the idea of an absolute man or woman is far from easy to grasp, not to mention that it cuts against the general commitment to egalitarianism in the modern world.

For the few interested in that path, two questions arise:

\begin{enumerate}


\item How do I know my nature?

\item How do I actualize it?
\end{enumerate}


These are intimately related since they represent knowing and being, which amount to the same thing. Esoteric psychology is the study of the soul, not empirically, but interiorly through direct observation. The more one knows about himself, the more being he will have. (The actual details will not be discussed here.)

Before proposing what the results may be, there are other competing notions for the ideal man or woman.

\paragraph{Better Living Through Chemistry}


In researching a recent post, I listened to a business consultant interviewed by a popular radio personality. The consultant's point of view was that actions should be taken that align with the ``hard-wiring'' of the brain. That is, the brain will respond by producing certain chemicals, specifically oxytocin, that makes us ``feel good.''

The interviewer agreed, and spoke about how good he felt following a demonstration, because people picked up the trash after themselves. Now, I'll agree that it is a moral good not to litter, but not because it produces certain brain chemicals. I'll remind readers about the serial killer, Ted Bundy, who told one of his victims how good it would feel to kill her.

\paragraph{Sexual Attraction}


The actor John Cusack recently complained that in Hollywood, an actress is over the hill at 29. That assumes that American film is an art form, in which acting ability is foremost. Actually, the women in most films are there to represent sexual attraction, and the age range of about 22 to 29 is the peak of that, for obvious reasons. So this sounds like a futile path for the absolute woman.

I recall years ago as a boy, that a nun told us that in heaven (or perhaps at the resurrection?) we would have bodies around the age 25, since that is the best age, physically speaking. I have never read that anywhere else, but I doubt she just made it up.

\paragraph{The Dalai Lama}


In earlier times, I had a rather romantic attachment to Tibet. Tibet was patriarchal, hierarchical, theocratic, autocratic, ancient, and traditional. There was occult knowledge and magical practices. There was the literature of travels through Tibet, the mythical Shambhala, and telepathic gurus.

My initiation and studies were very strict. I met many people who boasted to me that they had overcome their Christian religious upbringing. However, the requirements for the sangha were just as strict: no killing, no stealing, no lying, and no unnatural sex acts, as they are called in the West.

So it seems to me that the Fourteenth Dalai Lama has turned himself into a caricature of a holy man. He rejects traditional Tibet and would convert it to a Euro-style parliament system, even communistic, which he considers morally superior. He supports all the correct causes of the Left. It is difficult to see what Buddhism has to do with it.

Has he ever told an audience that their cravings cause their problems, they will likely be reborn as an insect, and their pet dogs are heading for the hell realms? I'm sure he never read \emph{The Doctrine of Awakening}.

\paragraph{George Clooney}


For many people, and for almost all women, George Clooney represents the ideal man. Handsome, suave, debonair, popular, wealthy, and articulate, he and the Dalai Lama support the same causes. However, Mr. Clooney is not at all religious, so his views must have come to him by some other means. If being fashionably liberal is your goal, just follow Mr. Clooney and dispense with the years of Buddhist meditation practice.

Mr. Clooney is the cosmopolitan man of today. He owns homes in the USA, Mexico, and Italy. He recently married a British woman from a prominent Lebanese family and the wedding was held in Rome. Her father is a Druze and her mother a Sunni Muslim. Just as Mr. Clooney is sanctified by his desire to end world hunger, the new Mrs. Clooney has her own high credentials. She is not an ambulance chaser, nor a sleazy corporate lawyer, but rather a ``human rights'' attorney.

Now, in times past, the idea of marrying an actor would have been considered a low-brow affair, but the ideals of feminine hypergamy have apparently changed.

\paragraph{The Warrior Ideal}


In the olden days, the warrior, particularly if he was also a political leader, was glorified. This was embodied in the ideal of the nine worthies\footnote{\url{https://gornahoor.net/?p=331}}. There are no current figures that approach that ideal, and more recent candidates are now too outré to extend the list to include three modern worthies. Perhaps Joan of Arc can fit in somehow. The worthies represent the state of the fallen world, in which it is necessary to toil and do battle against evil.

\paragraph{The Spiritual Ideal}


The highest European ideals before the modern era are certainly the Virgin Mary and Jesus. Without intending to deny more popular notions, from an esoteric point of view, they represent the two ideals we have already discussed.

\begin{itemize}


\item The Virgin Mary, free from sin, is our model of the Primordial State

\item Jesus, as true God and true man, represents \emph{theosis}, i.e., the raising of the human nature to the divine.
\end{itemize}


Now whereas those two figures were born that way, we have to actualize those states. Along with the worthies, these are the three possibilities open to us, as best suit our own natures.


\flrightit{Posted on 2014-10-03 by Cologero}

\begin{center}* * *\end{center}

\begin{footnotesize}
\begin{sffamily}
\texttt{Avery on 2014-10-03 at 10:18 said:}

You won't be seeing any great warriors glorified in the news, because we don't have any to report on. We don't even have American pilots any more, bursting through the sky in amazing machines going six thousand miles and hour. The American myth of the renegade fighter pilot is found in many of our greatest films. But now our pilots just sit in air conditioned rooms in Kansas and direct computers on a screen.

ISIS magazines are full of warriors; in one sense they are savvy traditionalists, while in another sense trying to immanentize the Eschaton.

\hfill

\texttt{Clayton on 2014-10-03 at 10:19 said:}

So now you are superior to the Dalai Lama. So your supposed Gnosis sure has not led to any inflation of your own self worth.

\hfill

\texttt{Cologero on 2014-10-03 at 12:44 said:}

If you read carefully, Clayton, I was referring to his social views which are totally Western and not representative of--and actually contrary to--any Tibetan tradition. Furthermore, those views owe nothing to the Buddhist tradition, since they are perfectly acceptable to those without any spiritual orientation at all. As an example of a contrary teaching, I pointed to sodomy, which, in his lineage, causes you to ``lose the vow from the root'', so you need a new initiation; he seems to have backed off from that. He also plans to dismantle the Tibetan social system if he gets the chance, he wants to abolish the dalai lama chain, and he plans to reincarnate as an attractive woman.

Then I mentioned some Tibetan teachings about which he is pretty much silent. Perhaps he has some esoteric reason for this. The point is, Clayton, that if you want to be an interesting and informed disputant, you would have addressed these \textbf{specific} points. None of those points you mentioned appears in the text. But thanks for trying.

\hfill

\texttt{rhondda9 on 2014-10-03 at 17:14 said:}

I have a soft spot for the Dali Lama. He seemed to exude joy. However, one day when a reporter asked him what he thought of all the westerners getting into Buddhism, he said that he thought they should look to their own traditions. I was quite taken aback and then realized how little I knew about Western Traditions. It was the beginning of the way home.

\hfill

\texttt{Clayton on 2014-10-03 at 18:50 said:}

I agree with you that he has moved away from some of these ``traditional'' notions which according to you automatically reduces him to a ``caricature of a Holy Man'' so apparently only those values which you hold in high esteem entitle one to Holiness. I see how you can disagree with a person's political views, but one's political views are in no way indicative of the degree of spirituality or holiness one has attained.

\hfill

\texttt{Cologero on 2014-10-03 at 19:11 said:}

You can continue to be mean-spirited, Clayton, but it is clouding your judgment. Who says I disagree with his politics? I, too, am for world peace, the end to hunger, a clean environment, and economic justice (although the Fourteenth Dalai Lama prefers communism). But holding such views does not require any particular spiritual transformation. If you were a deeper thinker, Clayton, you would have addressed the core issue. Why has the Dalai Lama abandoned traditional notions and adopted Western values in their stead? If Western values are superior, then who needs him? Or is he just playing a role that Westerners expect of a holy man, while keeping deeper teachings hidden?

\hfill

\texttt{Ash on 2014-10-04 at 00:23 said:}

I've just been re-reading Ride the Tiger. Evola discusses exactly this, the problem of having what are essentially multiple different personalities within the psyche. I think this multiplicity of tendencies is something I've always known about myself, but my introspection at this time is focusing on identifying my own ``central'' tendency. The actor as a symbol of the human in their most chaotic state is something completely counter to most people today. Funnily enough, I was just watching something this evening about Reagan, and they mentioned that his acting career was counted against him during his run for governor.

Avery: Strelkov, the Russian military adventurer, seems to be gaining warrior-hero status in his country. Whatever the true nature of the current Russian state, he does seem to believe that he is fighting for an imperial idea.

\hfill

\texttt{August on 2014-10-07 at 09:59 said:}

It's certainly hard to know what to make of those like the Dalai Lama when one learns some of their views... Perhaps they follow the old precept of not disturbing the laity, avoiding public discussion of ideas that go against the current of the times, and making do within the narrow confines of modern discourse, maybe working invisibly somehow. Since I cannot know for sure, I'll be gracious.

In other news, one of your ebola sufferers over there is supposedly a `tulku', one Ashoka Mukpo (see Wikipedia).

\end{sffamily}
\end{footnotesize}

